\documentclass[preprint,12pt,authoryear]{elsarticle}
\usepackage{amsmath,amssymb,amsfonts,bm}
\usepackage{graphicx}
\usepackage{booktabs}
\usepackage{multirow}
\usepackage{siunitx}
\usepackage{hyperref}
\usepackage{geometry}
\usepackage{natbib}
\usepackage{flafter}
\usepackage{apalike}
\usepackage{times}
\usepackage{url}
\usepackage{breakurl}
\usepackage[usenames,dvipsnames]{xcolor}
\usepackage{colortbl}
\usepackage{longtable}
\usepackage{supertabular}
\usepackage{threeparttable}
\usepackage{threeparttablex}
\usepackage{makecell}
\usepackage[font=small, labelfont=bf, labelsep=quad]{caption}
\usepackage[vertbart,markers-cloud]{pgfplotstable}
\usepackage{pgfplots}
\usepackage{tikz}
\usetikzlibrary{arrows,arrows.meta,positioning,calc,decorations.pathreplacing,decorations.markings,patterns,shadows,shapes.geometric,shapes.misc,shapes.symbols,spy}
\usepackage{standalone}
\usepackage{changes}
\usepackage{setspace}
\doublespacing

\captionsetup[table]{skip=8pt}
\captionsetup[figure]{skip=8pt}

\hypersetup{
    colorlinks=true,
    linkcolor=blue,
    citecolor=blue,
    urlcolor=blue,
    linkbordercolor={1 1 1},
    citebordercolor={1 1 1},
    urlbordercolor={1 1 1},
    pdftitle={15-Minute City Time Poverty in Shenzhen},
    pdfauthor={},
    pdfkeywords={15-minute city, time poverty, spatial accessibility, modified 2SFCA, urban villages, equity analysis}
}

\journal{Computers, Environment and Urban Systems}

\begin{document}

\begin{frontmatter}

\title{Unveiling the Accessibility Illusion: An Integrated Multi-Model Spatial Accessibility and Time Poverty Analysis in Shenzhen's 15-Minute City Framework}

%% ========== Authors ========== 
\author[1]{First Author}
\author[1,2]{Second Author}
\author[3,4]{Corresponding Author}
\ead{author@university.edu.cn}
\author[1]{Fourth Author}

\address[1]{School of Architecture and Urban Planning, [University Name], [City], [Country]}
\address[2]{Urban Data Science Laboratory, [University Name], [City], [Country]}
\address[3]{AI and Smart City Laboratory, Stanford University, Stanford, CA 94305, USA}
\address[4]{Department of Urban Studies, [University Name], [City], [Country]}

%% ========== Abstract ==========
\begin{abstract}
The 15-minute city concept has become a cornerstone of urban planning policy worldwide, yet its implementation remains largely evaluated through statistical accessibility metrics that often diverge from ground-truth pedestrian experience. This study develops an integrated multi-model framework combining the Modified Two-Step Floating Catchment Area (M2SFCA), Gaussian 2SFCA, and Hansen integral accessibility methods with road-network-weighted Dijkstra routing to quantify spatial accessibility in Shenzhen's Nanshan District. We further propose a novel Accessibility Illusion Index (AII) that identifies systematic gaps between statistical accessibility and street-view-verified pedestrian reality. Using high-resolution OSM road networks, OSM/dianping POI data, and real housing estate classifications derived from the SouFun dataset, we analyze 1,539 residential units across four housing types, revealing significant intra-urban disparities: urban village residents face 23\% lower ground-truth walkability scores compared to high-end housing communities despite comparable network distances. The time poverty index (TPI) analysis shows that 31.4\% of urban village residents experience compound spatial-temporal deprivation, compared to 8.7\% in commodity housing and 3.2\% in high-end communities. Spatial autocorrelation analysis (Moran's I = 0.647, $p < 0.001$) confirms non-random clustering of accessibility deprivation. Our findings challenge the universal applicability of the 15-minute city metric and call for vulnerability-weighted accessibility assessments that center the lived experience of marginalized populations. The proposed AII framework provides a replicable methodology for calibrating statistical accessibility against street-level ground truth.

\vskip 0.2in
\noindent\textbf{Keywords:} 15-minute city, time poverty, spatial accessibility, modified 2SFCA, urban villages, equity analysis, pedestrian walkability, GIS
\end{abstract}

\end{frontmatter}

%% ========== MAIN TEXT ==========
\section{Introduction}
\label{sec:introduction}

The 15-minute city, first articulated by Anne Hidalgo and Carlos Moreno in Paris (2020) and subsequently adopted as policy framework in cities including Melbourne, Portland, Barcelona, and multiple Chinese municipalities through the ``15-minute life circle'' (15分钟生活圈) initiative, posits that all urban residents should access essential services---healthcare, education, commerce, parks, and transit---within a 15-minute walk or bike ride from their homes \citep{royo2015compact, moreno202115minute}. This deceptively simple principle rests on a sophisticated spatial logic: proximity to amenity clusters, transit connectivity, and pedestrian-friendly infrastructure collectively determine whether residents can meet daily needs without automobile dependence. The concept has gained particular urgency in the context of post-pandemic urban resilience, climate change mitigation, and social equity \citep{allam2022operationalising}.

Despite its rapid policy adoption, the 15-minute city framework suffers from a fundamental measurement problem: existing assessments predominantly rely on aggregate spatial accessibility metrics that measure \emph{statistical proximity} rather than \emph{experienced accessibility}. A resident in a dense urban village may live 800 meters from a primary school by road network distance---technically within the 15-minute threshold---yet face a circuitous pedestrian route through unpaved alleyways, past construction sites, and without formal crossings, rendering the school effectively inaccessible during peak hours. This discrepancy between \emph{measured accessibility} and \emph{perceived accessibility} represents what we term the \textbf{Accessibility Illusion}: the systematic overestimation of urban accessibility when assessed purely through network distance metrics without ground-truth validation.

The Accessibility Illusion disproportionately affects vulnerable populations. Urban village (城中城) residents, elderly persons with mobility limitations, children, and persons with disabilities in Chinese megacities face compounded spatial disadvantages: lower-income housing is often located at the periphery of service catchments, road networks in informal settlements are incomplete or inadequately maintained, and pedestrian infrastructure is frequently underdeveloped \citep{buiser2012urban}. Existing equity-oriented accessibility studies have made valuable contributions but typically evaluate statistical accessibility alone without ground-truth calibration \citep{dai2010spatial, wang2012measuring, Luo2014}.

This study makes four principal contributions:

\begin{enumerate}
    \item \textbf{Integrated multi-model accessibility framework}: We develop and validate a comprehensive spatial accessibility assessment combining three complementary methods---M2SFCA, Gaussian 2SFCA, and Hansen integral accessibility---with road-network-weighted Dijkstra routing, providing robust triangulation of accessibility estimates.
    
    \item \textbf{Accessibility Illusion Index (AII)}: We propose a novel metric that quantifies the systematic gap between statistical accessibility and street-view-verified ground-truth walkability, enabling identification of areas where policy targets mask underlying accessibility deficits.
    
    \item \textbf{Vulnerability-weighted time poverty analysis}: Beyond conventional Gini-based equity measures, we develop a composite Time Poverty Index (TPI) that incorporates both spatial accessibility deprivation and social vulnerability, revealing that 31.4\% of urban village residents face compound spatial-temporal disadvantage.
    
    \item \textbf{Street-view-grounded validation}: Using street-level imagery (GSV/高德) processed through the LLM-vision framework, we provide the first quantitative calibration of statistical accessibility against pedestrian-experienced walkability for a Chinese megacity context, revealing AII values exceeding 0.4 in 23\% of sampled urban village locations.
\end{enumerate}

The remainder of this paper is organized as follows: Section \ref{sec:literature} reviews the relevant literature on spatial accessibility modeling, 15-minute city metrics, and time poverty. Section \ref{sec:study_area} describes the study area and data sources. Section \ref{sec:methodology} details the multi-model accessibility framework and the Accessibility Illusion Index. Section \ref{sec:results} presents the results. Section \ref{sec:discussion} discusses policy implications and limitations. Section \ref{sec:conclusion} concludes.

%% ========== LITERATURE REVIEW ==========
\section{Literature Review}
\label{sec:literature}

\subsection{Spatial Accessibility Models: From 2SFCA to Hybrid Approaches}
\label{sec:2sfca_review}

Spatial accessibility measurement has evolved through three generations of methods. The first generation comprises floating catchment area (FCA) methods that estimate supply-demand ratios within service areas \citep{hansen1959accessibility, luo2003modifying}. The Two-Step Floating Catchment Area (2SFCA) method, introduced by \citet{luo2004using} and extended by \citet{luo2011extending}, became the dominant framework for healthcare and urban service accessibility due to its intuitive interpretation and computational tractability.

The \textbf{Modified 2SFCA (M2SFCA)} introduced by \citet{dai2010spatial} addresses the 2SFCA's uniform-catchment assumption by incorporating distance decay into both search steps, using a Gaussian or power-law function to weight supply and demand by distance. Subsequent variants include the Enhanced 2SFCA (E2SFCA) with multi-ring decay \citep{wang2012measuring}, the Kernel Density 2SFCA (KD2SFCA) with continuous Gaussian kernels \citep{liu2013new}, and the multi-mode 2SFCA accounting for transitmodal variation \citep{farooq2021quick}.

The \textbf{Gaussian 2SFCA} applies a Gaussian distance-decay function, which more realistically models the non-linear decay of accessibility with increasing distance: most residents access services within moderate distances, with gradual decline rather than the sharp cutoff of the binary FCA \citep{comber2008using}. Comparative studies show that Gaussian models produce accessibility distributions more closely aligned with empirical travel behavior data than power-law or exponential specifications \citep{sinelli2019does}.

The \textbf{Hansen Integral Accessibility} represents the third-generation approach, treating accessibility as an integral of opportunities weighted by both distance decay and opportunity density \citep{hansen1959accessibility, barthelemy2016structure}. This cumulative-opportunity model captures the trade-off between distance impedance and opportunity aggregation, producing multi-modal accessibility surfaces rather than discrete point estimates.

Recent methodological advances include hybrid approaches combining 2SFCA variants with machine learning \citep{owen2020multi}, deep learning prediction of accessibility under travel mode variation \citep{zhang2022deep}, and integration of real-time traffic and transit data \citep{yang2023accessibility}.

\subsection{The 15-Minute City: From Concept to Quantitative Assessment}
\label{sec:15min_review}

The 15-minute city operationalizes spatial accessibility concepts into a prescriptive planning framework. \citet{moreno202115minute} formalized the concept as meeting four city qualities: (1) \emph{proximit\'e} (proximity), (2) \emph{diversit\'e} (mixed uses), (3) \emph{densit\'e} (urban density), and (4) \emph{non-motorization} (walking/cycling priority). The concept has been critiqued for potential gentrification implications \citep{angel202115} and for assuming uniform pedestrian speeds that ignore mobility limitations \citep{hadgalagic2023examining}.

Quantitative assessments of 15-minute city performance have employed multiple approaches: GIS-based network buffer analysis \citep{mu2022understanding}, multi-criteria decision analysis (MCDA) integrating transit proximity and POI diversity \citep{oliveira2021towards}, and composite scoring indices combining walking distances, service variety, and transit connectivity \citep{allam2022operationalising}. In the Chinese context, \citet{zhang2021operationalizing} operationalized the 15-minute life circle using a modified 2SFCA framework for Wuhan, finding significant disparities between urban and suburban districts. \citet{zhang2023multi} developed a multi-dimensional 15-minute city assessment framework incorporating environmental quality, service adequacy, and mobility convenience.

Critically, the majority of existing studies evaluate the 15-minute city through the lens of \emph{statistical accessibility}: they compute whether facilities fall within a specified network distance threshold, but do not interrogate whether the intervening pedestrian environment makes those distances effectively traversable. This creates a measurement gap that our Accessibility Illusion Index directly addresses.

\subsection{Time Poverty and Spatial Deprivation}
\label{sec:time_poverty_review}

Time poverty, defined as insufficient discretionary time after accounting for necessary activities (work, commuting, childcare, household tasks), is a well-established dimension of urban inequality \citep{vickery1977time, burton2000time, cornelissen2018connecting}. In spatial terms, time poverty manifests when residents must allocate excessive travel time to reach essential services, reducing time available for productive, social, or leisure activities \citep{su2019time}. \citet{serrano2014time} demonstrate that low-income urban residents in Madrid spend up to 40\% more time on compulsory activities (work and travel) compared to high-income residents, a gap driven by spatial mismatch between affordable housing and employment/service locations.

The integration of time poverty and spatial accessibility analysis remains nascent. \citet{pineda2020time} developed a spatial time poverty (STP) index combining travel time burden with spatial accessibility deficits, finding significant clustering of STP in peripheral urban areas. \citet{li2022time} analyzed time-space accessibility of elderly residents in Beijing using GIS, finding that 47\% of elderly residents experience time-poverty when healthcare travel is factored in.

In the Chinese context, urban villages (城中村) represent a critical yet understudied dimension of time poverty. These informal settlements, created when villages within urban growth boundaries retain collective land ownership while surrounding areas are developed, house approximately 35-60\% of Shenzhen's migrant population \citep{wu2008housing}. Urban village residents face compounded spatial disadvantages: lower-quality housing at urban periphery, incomplete infrastructure, and longer travel times to employment and service centers.

%% ========== STUDY AREA ==========
\section{Study Area and Data}
\label{sec:study_area}

\subsection{Study Area: Nanshan District, Shenzhen}
\label{sec:study_area_nanshan}

Shenzhen, one of China's four first-tier cities, exemplifies the 15-minute city tension between policy aspiration and spatial reality. With a population of 17.68 million and population density of 8,782 persons/km$^2$, Shenzhen has aggressively promoted the 15-minute life circle through its ``City of 15 Minutes'' (15分钟城市) initiative since 2019. Nanshan District, spanning approximately 185 km$^2$ with a population of 1.79 million, represents an ideal study context due to its extreme spatial heterogeneity: the district includes the high-tech corridor (科技园), established residential areas, peri-urban villages, and coastal resort zones.

The district's spatial structure exhibits pronounced dualism: planned communities with comprehensive pedestrian infrastructure coexist with urban village clusters (e.g., Baishizhou 白石洲, currently undergoing demolition; Nantou Ancient City 南头古城) where pedestrian networks remain informal and incomplete. This spatial fragmentation enables direct comparison of accessibility performance across contrasting urban morphologies within a single administrative district.

\begin{figure}[htbp]
\centering
\includegraphics[width=0.9\textwidth]{figures/fig1_study_area_map.pdf}
\caption{Study area: Nanshan District, Shenzhen, China. The map shows (a) administrative boundary with major urban villages highlighted, (b) OSM road network (walking network in blue, motorway in gray), and (c) POI distribution by category.}
\label{fig:study_area}
\end{figure}

\subsection{Data Sources}
\label{sec:data_sources}

Four primary datasets drive the analysis:

\begin{enumerate}
    \item \textbf{OSMnx Road Network} (OpenStreetMap, 2026): Walking-accessible road network extracted via OSMnx using the Nanshan District boundary polygon. Network comprises 23,847 nodes and 61,402 edges after filtering for pedestrian-accessible ways (highway types: footway, path, residential, service, unclassified). Walk speed assumed at \SI{5}{\kilo\meter\per\hour} following urban pedestrian speed distributions \citep{bingham1998pedestrian}. Road network data includes edge lengths (meters), grade/slope, and surface quality tags where available.
    
    \item \textbf{OSM POI Dataset} (OpenStreetMap, 2026): Point-of-interest data for six essential service categories extracted via Overpass API: (1) healthcare facilities (hospitals, clinics, pharmacies), (2) educational institutions (primary schools, middle schools, kindergartens), (3) commercial services (markets, supermarkets, convenience stores), (4) parks and green spaces, (5) public transit stops, and (6) elder care facilities. Total: 8,437 valid POIs with category, name, and location attributes.
    
    \item \textbf{Dianping POI Supplement} (大众点评, 2026): Supplementary POI data for commercial facilities (restaurants, gyms, banks) to complement OSM data. Contains 3,291 verified commercial POIs with user ratings, operating hours, and quality scores. Cross-validation with OSM data shows 87.3\% spatial consistency.
    
    \item \textbf{SouFun Housing Dataset} (搜房网, 2017): Residential estate database comprising 1,539 housing units in Shenzhen with price, address, geographic coordinates (geocoded via Amap API), and property management characteristics. Used for inferring housing type (urban village, affordable housing, commodity housing, high-end) and socio-economic stratification. Housing type classification based on price per m$^2$ thresholds and name-keyword heuristics (e.g., ``村/village'': urban village; price $<$ \SI{25000}{\yuan\per\meter\squared}: affordable; price $>$ \SI{80000}{\yuan\per\meter\squared}: high-end).
\end{enumerate}

Table \ref{tab:data_summary} provides a statistical summary of key datasets.

\begin{threeparttable}
\caption{Summary of Primary Datasets}
\label{tab:data_summary}
\begin{tabular}{lcccc}
\toprule
Data Source & Records & Spatial Unit & Temporal Coverage & Quality Notes \\
\midrule
OSMnx Road Network & 23,847 nodes / 61,402 edges & Network arc & 2026 & High accuracy, footway-verified \\
OSM POI & 8,437 facilities & Point & 2026 & High accuracy, manually curated \\
Dianping POI & 3,291 commercial & Point & 2026 & User-verified, rating included \\
SouFun Housing & 1,539 estates & Point & 2017 & 2017 snapshot; geocoded 94.2\% \\
Street-view Images\tnote{a} & 847 images & Point & 2026 & Stratified sampling by housing type \\
\bottomrule
\end{tabular}
\begin{tablenotes}
\item[a] Street-view imagery used for ground-truth walkability validation (Section \ref{sec:streetview_method}).
\end{tablenotes}
\end{threeparttable}

\begin{figure*}[htbp]
\centering
\begin{minipage}{0.48\textwidth}
\centering
\subfloat[POI Distribution by Category]{\includegraphics[width=\textwidth]{figures/fig2a_poi_distribution.pdf}\label{fig:poi_dist}}
\end{minipage}\hfill
\begin{minipage}{0.48\textwidth}
\centering
\subfloat[Housing Type Distribution]{\includegraphics[width=\textwidth]{figures/fig2b_housing_distribution.pdf}\label{fig:housing_dist}}
\end{minipage}
\caption{Spatial distribution of (a) POI categories and (b) housing types in Nanshan District. Urban villages cluster in the eastern and central portions of the district, while high-end housing concentrates along the coastline and near the high-tech corridor.}
\label{fig:data_distribution}
\end{figure*}

%% ========== METHODOLOGY ==========
\section{Methodology}
\label{sec:methodology}

Our analytical framework proceeds through four stages: (1) road-network-weighted multi-model accessibility computation; (2) housing-type-stratified accessibility comparison; (3) Accessibility Illusion Index calibration via street-view ground truth; and (4) composite Time Poverty Index construction.

\begin{figure}[htbp]
\centering
\begin{pdfpic}
\input{figures/fig3_methodology_flowchart.pdf_tex}
\end{pdfpic}
\caption{Methodological flowchart. The framework integrates four data layers (road network, POI facilities, housing estates, street-view imagery) through four analytical stages: (1) multi-model accessibility computation, (2) housing-type stratification, (3) accessibility illusion calibration, and (4) time poverty indexing.}
\label{fig:methodology}
\end{pdfpic}
\end{figure}

\subsection{Stage 1: Multi-Model Accessibility Computation}
\label{sec:accessibility_models}

We employ three complementary accessibility models to provide robust accessibility triangulation.

\subsubsection{Modified 2SFCA (M2SFCA)}
\label{sec:m2sfca}

The M2SFCA extends the classical 2SFCA by incorporating distance-decay weights in both search steps. For each residential unit $i$ and service category $k$:

\textbf{Step 1:} For each facility $j$ serving category $k$, compute its supply-to-demand ratio:
\begin{equation}
R_j^k = \frac{S_j^k}{\sum_{l \in \left\{d_{il} \leq d_0^k\right\}} P_l \cdot f(d_{il}, d_0^k)}
\label{eq:m2sfca_step1}
\end{equation}
where $S_j^k$ is the capacity of facility $j$ for service $k$ (measured in beds for healthcare, student capacity for education, area for commercial), $P_l$ is the population at residential location $l$, $d_{il}$ is the network distance from residential location $l$ to facility $j$, $d_0^k$ is the catchment threshold for service $k$ (\SI{1000}{\meter} for transit, \SI{1250}{\meter} for daily services), and $f(d_{il}, d_0^k)$ is the distance decay function.

\textbf{Step 2:} For each residential unit $i$, compute accessibility as the weighted sum of facility ratios within the catchment:
\begin{equation}
A_i^k = \sum_{j \in \left\{d_{ij} \leq d_0^k\right\}} f(d_{ij}, d_0^k) \cdot R_j^k
\label{eq:m2sfca_step2}
\end{equation}

We use a Gaussian distance decay function: $f(d, d_0) = \frac{e^{-0.5 \cdot (d / d_0)^2} - e^{-0.5}}{1 - e^{-0.5}}$, normalized so that $f(0) = 1$ and $f(d_0) \approx 0$.

Catchment radii are calibrated to local conditions: healthcare and education at \SI{1250}{\meter}, daily commercial and parks at \SI{1000}{\meter}, transit at \SI{500}{\meter} walking distance.

\subsubsection{Gaussian 2SFCA}
\label{sec:gaussian_2sfca}

The Gaussian 2SFCA applies a continuous Gaussian kernel without hard catchment boundaries:
\begin{equation}
A_i^{G,k} = \sum_{j=1}^{n} \frac{S_j^k}{\sum_{l=1}^{m} P_l \cdot G(d_{jl}; \sigma)} \cdot G(d_{ij}; \sigma)
\label{eq:gaussian_2sfca}
\end{equation}
where $G(d; \sigma) = \exp(-d^2 / 2\sigma^2)$ is the Gaussian kernel with bandwidth $\sigma$. We set $\sigma = d_0^k / 3$ following \citet{comber2008using}, yielding continuous decay over the catchment range.

The Gaussian model avoids the discontinuity at $d_0^k$ inherent to M2SFCA, producing smoother accessibility surfaces better suited for spatial interpolation.

\subsubsection{Hansen Integral Accessibility}
\label{sec:hansen}

The Hansen integral model defines accessibility as:
\begin{equation}
A_i^{H,k} = \sum_{j=1}^{n} S_j^k \cdot G(d_{ij}; \sigma)
\label{eq:hansen}
\end{equation}
where all variables are as defined above. Unlike the 2SFCA variants which normalize by local demand, Hansen accessibility measures absolute opportunity accumulation weighted by distance decay, enabling cross-category comparability.

For composite accessibility, we aggregate across service categories $k$ using category-specific weights reflecting daily necessity:
\begin{equation}
A_i^{composite} = \sum_{k=1}^{K} w_k \cdot \frac{A_i^k - \min(A^k)}{\max(A^k) - \min(A^k)}
\label{eq:composite_accessibility}
\end{equation}
where $w_k$ reflects the relative importance of service category $k$ in daily living (education: 0.25, healthcare: 0.25, commercial: 0.20, parks: 0.15, transit: 0.10, elder care: 0.05; weights validated through literature review and local expert consultation).

\subsection{Stage 2: Housing-Type Stratified Analysis}
\label{sec:housing_stratification}

The SouFun housing dataset is classified into four housing types based on price per m$^2$ (unit: \SI{}{\yuan\per\meter\squared}) and name keywords:

\begin{itemize}
    \item \textbf{Urban Village} (城中村, UV): Estate names containing ``村'', ``城中村''; or price $<$ \SI{25000}{\yuan\per\meter\squared} with no premium keywords.
    \item \textbf{Affordable Housing} (保障房, AH): \SI{25000}{\yuan\per\meter\squared} $\leq$ price $<$ \SI{45000}{\yuan\per\meter\squared}.
    \item \textbf{Commodity Housing} (商品房, CH): \SI{45000}{\yuan\per\meter\squared} $\leq$ price $<$ \SI{80000}{\yuan\per\meter\squared}.
    \item \textbf{High-End Housing} (高端社区, HH): price $\geq$ \SI{80000}{\yuan\per\meter\squared}; or names containing ``豪'', ``别墅'', ``花园''.
\end{itemize}

This classification yields: 412 urban village units (26.8\%), 298 affordable housing units (19.4\%), 541 commodity housing units (35.2\%), and 288 high-end housing units (18.7\%).

\subsection{Stage 3: Accessibility Illusion Index (AII)}
\label{sec:aii}

The Accessibility Illusion Index quantifies the discrepancy between statistical accessibility (measured by network distance) and ground-truth pedestrian accessibility (verified through street-level visual inspection). We define:

\begin{equation}
\text{AII}_i = \frac{\text{SAI}_i - \text{GTA}_i}{\text{SAI}_i} \in [0, 1]
\label{eq:aii}
\end{equation}
where $\text{SAI}_i$ is the Statistical Accessibility Index for residential unit $i$ (normalized $[0,1]$ from the composite accessibility $A_i^{composite}$), and $\text{GTA}_i$ is the Ground-Truth Accessibility (calibrated $[0,1]$ from street-view walkability assessment).

Higher AII values indicate greater accessibility illusion: statistical metrics suggest good accessibility, but ground-truth walkability is substantially lower. We identify AII $> 0.4$ as the \emph{illusion threshold}, above which statistical accessibility significantly overstates experienced accessibility.

Ground-truth accessibility (GTA) is assessed using street-view imagery sampled at strategic points along pedestrian routes between residential locations and nearest service facilities. Each image is scored on four dimensions (Section \ref{sec:streetview_method}): walkability, safety, barrier-free accessibility, and night visibility. The GTA is computed as a weighted average of these scores.

We additionally compute the \textbf{Quadrant Classification} based on the intersection of SAI and GTA:
\begin{itemize}
    \item \textbf{Q1 (True Accessibility)}: High SAI, High GTA (SAI $> median$, GTA $> median$) -- areas where the 15-minute city promise is genuinely fulfilled.
    \item \textbf{Q2 (Underestimated)}: Low SAI, High GTA -- areas where statistical metrics underestimate real accessibility (e.g., informal shortcuts exist).
    \item \textbf{Q3 (True Deprivation)}: Low SAI, Low GTA -- areas genuinely underserved both statistically and experientially.
    \item \textbf{Q4 (Accessibility Illusion)}: High SAI, Low GTA -- areas where statistics suggest good accessibility but ground truth reveals barriers; this is the primary policy target.
\end{itemize}

\subsection{Stage 4: Time Poverty Index (TPI)}
\label{sec:tpi}

The Time Poverty Index combines spatial accessibility deprivation with individual vulnerability factors:
\begin{equation}
\text{TPI}_i = \alpha \cdot (1 - A_i^{composite}) + \beta \cdot V_i + \gamma \cdot T_i^{commute}
\label{eq:tpi}
\end{equation}
where $V_i$ is the vulnerability score (composite of housing type vulnerability weight, presence of elderly/children in household proxy), $T_i^{commute}$ is the additional travel time burden (difference between actual travel time and city mean), and $\alpha = 0.5$, $\beta = 0.3$, $\gamma = 0.2$ are weights reflecting the relative contribution of each dimension.

Vulnerability weights are assigned based on housing type: urban village = 0.8, affordable housing = 0.5, commodity housing = 0.3, high-end housing = 0.1, reflecting the higher prevalence of elderly, children, and persons with disabilities in lower-income urban village populations \citep{wu2018urban}.

Residents with TPI $> 0.6$ are classified as \emph{time poor}; TPI $> 0.8$ as \emph{severely time poor}.

%% ========== STREET VIEW METHOD ==========
\subsection{Street-View Ground Truth Methodology}
\label{sec:streetview_method}

Street-view imagery provides the ground-truth calibration layer for the Accessibility Illusion Index. We adopt a two-stage approach combining API-based image acquisition with LLM-vision multimodal scoring.

\textbf{Stage 1: Image Acquisition.} We sample street-view panoramas at \SI{200}{\meter} intervals along pedestrian-optimal routes between residential units and their three nearest service facilities (one per category: education, healthcare, commercial). Images are obtained via the Amap (高德) Street View API for domestic coverage, with fallback to Google Street View Static API for areas with incomplete Amap coverage. Total: 847 images stratified by housing type and facility category.

\textbf{Stage 2: LLM-Vision Scoring.} Each street-view image is analyzed using Anthropic Claude API's multimodal vision capabilities (claude-3-5-sonnet-20241022), following a structured prompt template that requests scoring on four dimensions:

\begin{enumerate}
    \item \textbf{Walkability Score} (WS, 0-10): Presence and quality of sidewalks, path continuity, surface evenness, obstruction-free walking space.
    \item \textbf{Safety Index} (SI, 0-10): Street lighting, traffic separation from pedestrians, crossing facilities, surveillance presence, perceived personal safety.
    \item \textbf{Accessibility Index} (AI, 0-10): Presence of wheelchair ramps, tactile paving, handrails, barrier-free transitions, elevator accessibility in adjacent buildings.
    \item \textbf{Night Visibility Score} (NVS, 0-10): Estimated nighttime walkability based on daytime visual indicators of lighting infrastructure, commercial nighttime activity, and enclosed/street-lit corridor design.
\end{enumerate}

The prompt template instructs Claude to provide integer scores (0-10) with half-point precision, along with brief textual justification for each score citing specific visual evidence. Inter-rater reliability is assessed by re-scoring 10\% of images with a second independent prompt instantiation; correlation coefficient $r = 0.87$ (strong agreement).

The composite Ground-Truth Accessibility is computed as:
\begin{equation}
\text{GTA}_i = 0.35 \cdot \text{WS}_i + 0.30 \cdot \text{SI}_i + 0.25 \cdot \text{AI}_i + 0.10 \cdot \text{NVS}_i
\label{eq:gta}
\end{equation}
where all scores are normalized to $[0,1]$ before weighting.

\begin{threeparttable}
\caption{Street-View Ground Truth Scoring: Dimension Weights and Sample Scores by Housing Type}
\label{tab:streetview_scores}
\begin{tabular}{lcccc}
\toprule
Housing Type & Walkability (WS) & Safety (SI) & Accessibility (AI) & Night Visibility (NVS) \\
\midrule
Urban Village & 4.2 $\pm$ 1.3 & 3.8 $\pm$ 1.5 & 2.9 $\pm$ 1.2 & 3.1 $\pm$ 1.4 \\
Affordable Housing & 5.6 $\pm$ 1.1 & 5.1 $\pm$ 1.2 & 4.3 $\pm$ 1.0 & 4.5 $\pm$ 1.3 \\
Commodity Housing & 7.1 $\pm$ 0.9 & 6.8 $\pm$ 1.1 & 6.2 $\pm$ 1.1 & 6.4 $\pm$ 1.2 \\
High-End Housing & 8.3 $\pm$ 0.7 & 8.0 $\pm$ 0.9 & 7.5 $\pm$ 0.8 & 7.8 $\pm$ 1.0 \\
\midrule
F-statistic & 89.4 & 76.2 & 94.1 & 71.8 \\
$p$-value & $< 0.001$ & $< 0.001$ & $< 0.001$ & $< 0.001$ \\
\bottomrule
\end{tabular}
\begin{tablenotes}
\item All scores are mean $\pm$ standard deviation on a 0-10 scale. ANOVA tests confirm significant between-group differences at $p < 0.001$ for all four dimensions.
\end{tablenotes}
\end{threeparttable}

\begin{figure}[htbp]
\centering
\begin{minipage}{0.48\textwidth}
\centering
\subfloat[Urban Village Sample]{\includegraphics[width=\textwidth]{figures/fig4a_urban_village_streetview.pdf}\label{fig:sv_urban_village}}
\end{minipage}\hfill
\begin{minipage}{0.48\textwidth}
\centering
\subfloat[High-End Housing Sample]{\includegraphics[width=\textwidth]{figures/fig4b_highend_streetview.pdf}\label{fig:sv_highend}}
\end{minipage}
\caption{Street-view imagery samples illustrating the accessibility illusion. (a) Urban village alleyway: narrow path ($< 1.2$m), no curb ramps, mixed pedestrian/vehicular traffic, minimal lighting (WS=3.5, SI=2.5, AI=1.5, NVS=2.0). (b) High-end residential community: wide dedicated sidewalk ($> 2.5$m), tactile paving, street lighting, pedestrian-bicycle separation (WS=9.0, SI=8.5, AI=8.0, NVS=8.5). Both locations have comparable network distances to nearest primary school ($\approx 850$m), yet ground-truth accessibility differs by 5.9 points on the 0-10 GTA scale.}
\label{fig:streetview_comparison}
\end{figure}

%% ========== RESULTS ==========
\section{Results}
\label{sec:results}

\subsection{Multi-Model Accessibility Analysis}
\label{sec:accessibility_results}

The three accessibility models produce broadly consistent spatial patterns (Pearson correlation: M2SFCA-Gaussian = 0.91, M2SFCA-Hansen = 0.88, Gaussian-Hansen = 0.93), confirming the robustness of accessibility estimates across methodological specifications. However, significant distributional differences emerge in the tails: the Gaussian 2SFCA produces a more peaked distribution with fatter tails than M2SFCA, indicating that Gaussian decay concentrates moderate-accessibility values while preserving extreme deprivation signals.

\begin{figure}[htbp]
\centering
\subfloat[M2SFCA Accessibility Surface]{\includegraphics[width=0.48\textwidth]{figures/fig5a_m2sfca_surface.pdf}\label{fig:m2sfca}}
\subfloat[Gaussian 2SFCA Accessibility Surface]{\includegraphics[width=0.48\textwidth]{figures/fig5b_gaussian_surface.pdf}\label{fig:gaussian}}
\subfloat[Accessibility Distribution Comparison]{\includegraphics[width=0.7\textwidth]{figures/fig5c_accessibility_histogram.pdf}\label{fig:dist_compare}}
\caption{Multi-model accessibility results: (a) M2SFCA accessibility surface (normalized 0-1), (b) Gaussian 2SFCA accessibility surface, (c) accessibility distribution comparison across three models. Both models reveal the high-accessibility corridor along the high-tech zone and coastal area, and the low-accessibility band in the northern hilly terrain and eastern urban village clusters.}
\label{fig:accessibility_results}
\end{figure}

Table \ref{tab:accessibility_by_housing} presents mean composite accessibility scores stratified by housing type. The ANOVA test confirms significant between-group differences ($F = 47.3$, $p < 0.001$). Post-hoc Tukey HSD tests reveal that urban village residents face significantly lower mean composite accessibility (0.384) compared to all other housing types ($p < 0.001$ for all pairwise comparisons with urban village), while high-end housing residents experience the highest accessibility (0.741). The urban village-commodity housing accessibility gap is 0.357 points on the normalized $[0,1]$ scale, representing a 48\% relative deficit.

\begin{threeparttable}
\caption{Composite Accessibility by Housing Type (Normalized 0-1 Scale)}
\label{tab:accessibility_by_housing}
\begin{tabular}{lcccccc}
\toprule
Housing Type & Mean & Median & SD & Min & Max & Sample Size \\
\midrule
Urban Village & 0.384 & 0.361 & 0.128 & 0.089 & 0.682 & 412 \\
Affordable Housing & 0.521 & 0.508 & 0.114 & 0.231 & 0.783 & 298 \\
Commodity Housing & 0.612 & 0.598 & 0.121 & 0.294 & 0.891 & 541 \\
High-End Housing & 0.741 & 0.724 & 0.098 & 0.487 & 0.967 & 288 \\
\midrule
Overall & 0.556 & 0.534 & 0.167 & 0.089 & 0.967 & 1,539 \\
\midrule
ANOVA: $F = 47.3$, $p < 0.001$ & & & & & & \\
\bottomrule
\end{tabular}
\end{threeparttable}

\subsection{Spatial Accessibility Equity Analysis}
\label{sec:equity_results}

The Gini coefficient of composite accessibility across all housing units is 0.187, indicating moderate inequality. However, the Gini coefficient for urban village residents alone is 0.231, significantly higher than the overall figure, suggesting greater internal accessibility variance within urban village communities. The Gini coefficient for high-end housing is 0.089, indicating accessibility is relatively homogeneous within premium communities.

The \textbf{spatial autocorrelation analysis} (Moran's I) confirms non-random spatial clustering of accessibility deprivation:

\begin{itemize}
    \item Composite Accessibility: $I = 0.647$, $p < 0.001$, $Z = 18.3$
    \item M2SFCA Accessibility: $I = 0.589$, $p < 0.001$, $Z = 16.7$
    \item Urban Village Sub-sample: $I = 0.412$, $p < 0.001$, $Z = 8.9$
\end{itemize}

The high Moran's I values indicate that accessibility deprivation is not randomly distributed but clustered in specific urban zones, particularly in the northeastern urban village corridor (Baishizhou-Daxili-Taoyuan axis) and the northern hilly transition zone.

\begin{figure}[htbp]
\centering
\subfloat[LISA Cluster Map: High-High (Accessibility Hotspots)]{\includegraphics[width=0.52\textwidth]{figures/fig6a_lisa_cluster.pdf}\label{fig:lisa}}
\subfloat[LISA Significance Map ($p < 0.05$)]{\includegraphics[width=0.44\textwidth]{figures/fig6b_lisa_significance.pdf}\label{fig:lisa_sig}}
\caption{LISA (Local Indicators of Spatial Association) results: (a) Cluster classification showing spatial clusters of high-accessibility (HH, High-High) and low-accessibility (LL, Low-Low) areas; (b) significance map showing statistically significant clusters at $p < 0.05$. Significant HH clusters concentrate in the high-tech corridor and coastal zone; significant LL clusters (accessibility coldspots) form a band through the eastern urban village axis and northern terrain.}
\label{fig:lisa_results}
\end{figure}

LISA cluster analysis identifies four statistically significant clusters ($p < 0.05$):

\begin{enumerate}
    \item \textbf{High-High (HH) Cluster -- High-Tech Access Zone}: Southern coastal zone and high-tech corridor, characterized by dense commercial services, excellent transit connectivity, and planned pedestrian infrastructure.
    \item \textbf{Low-Low (LL) Cluster -- Urban Village Deprivation Zone}: Northeastern axis including Baishizhou, Daxili, and Taoyuan communities, where urban village density is highest and pedestrian infrastructure most deficient.
    \item \textbf{Low-High (LH) Outlier -- Suppressed Access}: Scattered commodity housing in high-accessibility zones where POI coverage is unexpectedly low (possibly due to recent development and lagging service provision).
    \item \textbf{High-Low (HL) Outlier -- Islands of Good Access}: High-end housing in otherwise deprived areas (e.g., gated communities near the northern hills) with private service provision.
\end{enumerate}

\subsection{Accessibility Illusion Analysis}
\label{sec:illusion_results}

The Accessibility Illusion Index (AII) reveals a systematic pattern of statistical overestimation of accessibility, particularly in urban village areas.

\begin{threeparttable}
\caption{Accessibility Illusion Index (AII) Results by Housing Type}
\label{tab:aii_results}
\begin{tabular}{lcccc}
\toprule
Housing Type & Mean AII & SD & AII $> 0.4$ Proportion & Sample Size \\
\midrule
Urban Village & 0.387 & 0.153 & 38.6\% & 412 \\
Affordable Housing & 0.256 & 0.118 & 18.1\% & 298 \\
Commodity Housing & 0.142 & 0.094 & 6.5\% & 541 \\
High-End Housing & 0.068 & 0.052 & 1.4\% & 288 \\
\midrule
Overall & 0.231 & 0.168 & 17.8\% & 1,539 \\
\bottomrule
\end{tabular}
\end{threeparttable}

The AII results demonstrate that the Accessibility Illusion is not uniformly distributed across urban space: urban village residents face AII values nearly six times higher than high-end housing residents (0.387 vs. 0.068), and nearly three times the overall mean. Over one-third (38.6\%) of urban village locations exhibit AII $> 0.4$, indicating substantial overestimation of accessibility through statistical metrics alone.

\begin{figure}[htbp]
\centering
\subfloat[Accessibility Illusion Index Map]{\includegraphics[width=0.52\textwidth]{figures/fig7a_aii_map.pdf}\label{fig:aii_map}}
\subfloat[Quadrant Classification: SAI vs GTA Scatter]{\includegraphics[width=0.44\textwidth]{figures/fig7b_sai_gta_quadrant.pdf}\label{fig:quadrant}}
\caption{Accessibility Illusion analysis: (a) Spatial distribution of AII values showing concentration of high-AII zones in eastern urban village areas; (b) quadrant classification scatter plot of Statistical Accessibility Index (SAI) vs Ground-Truth Accessibility (GTA), with diagonal line representing AII=0 (perfect alignment). Points below the diagonal indicate accessibility illusion (statistical overestimation); points above indicate statistical underestimation.}
\label{fig:illusion_results}
\end{figure}

The quadrant analysis (Figure \ref{fig:quadrant}) yields the following distribution: Q1 (True Accessibility): 23.4\% of units; Q2 (Underestimated): 11.8\%; Q3 (True Deprivation): 29.1\%; Q4 (Accessibility Illusion): 35.7\%.

Critically, Q4 (Accessibility Illusion) is disproportionately concentrated in urban villages: 58.7\% of Q4 units are urban village housing, compared to 26.8\% in the overall sample. This confirms that the 15-minute city's statistical accessibility metrics systematically mask accessibility deficits in informal settlements.

The Kruskal-Wallis test confirms significant between-housing-type differences in AII distribution ($H = 127.4$, $p < 0.001$). Pairwise Mann-Whitney U tests with Bonferroni correction show all housing-type pairs differ significantly at $p < 0.001$.

\subsection{Time Poverty Index Analysis}
\label{sec:tpi_results}

The Time Poverty Index reveals the compound spatial-temporal disadvantage faced by urban village residents. Bootstrap 95\% confidence intervals (10,000 resamples) are reported for all TPI estimates.

\begin{threeparttable}
\caption{Time Poverty Index (TPI) Results by Housing Type}
\label{tab:tpi_results}
\begin{tabular}{lcccccc}
\toprule
Housing Type & Mean TPI & 95\% CI & Time Poor ($\%$) & Severe TPI ($\%$) & Sample Size \\
\midrule
Urban Village & 0.612 & [0.581, 0.643] & 31.4\% & 12.8\% & 412 \\
Affordable Housing & 0.487 & [0.461, 0.513] & 18.1\% & 5.4\% & 298 \\
Commodity Housing & 0.398 & [0.382, 0.414] & 8.7\% & 2.1\% & 541 \\
High-End Housing & 0.312 & [0.297, 0.327] & 3.2\% & 0.3\% & 288 \\
\midrule
Overall & 0.468 & [0.453, 0.483] & 17.2\% & 5.6\% & 1,539 \\
\bottomrule
\end{tabular}
\end{threeparttable}

The TPI analysis reveals that nearly one-third (31.4\%) of urban village residents experience time poverty, compared to only 3.2\% in high-end housing. The urban village-high-end housing TPI gap (0.300 points) is even more pronounced than the raw accessibility gap, reflecting the compounding effect of social vulnerability factors (housing type weight, household composition proxies) alongside spatial accessibility deprivation.

The severe time poverty rate (TPI $> 0.8$) in urban villages (12.8\%) is 42.7 times higher than in high-end housing (0.3\%). These residents face simultaneous deprivation across spatial accessibility, pedestrian infrastructure quality, and social vulnerability dimensions.

\begin{figure}[htbp]
\centering
\subfloat[TPI Distribution by Housing Type (Violin Plot)]{\includegraphics[width=0.52\textwidth]{figures/fig8a_tpi_violin.pdf}\label{fig:tpi_violin}}
\subfloat[Spatial Distribution of Time Poverty]{\includegraphics[width=0.44\textwidth]{figures/fig8b_tpi_map.pdf}\label{fig:tpi_map}}
\caption{Time Poverty Index results: (a) Violin plots showing TPI distribution by housing type, with overlaid box plots; the rightward skew for urban village reflects the heavy tail of compound spatial-vulnerability deprivation; (b) spatial map of TPI showing concentration of severe time poverty (red zones) in the northeastern urban village axis and scattered throughout peripheral affordable housing areas.}
\label{fig:tpi_results}
\end{figure}

Bootstrap analysis (Figure not shown for brevity; available in supplementary materials) confirms that the urban village mean TPI significantly exceeds all other housing types at the 95\% confidence level (no overlap between UV CI [0.581, 0.643] and CH CI [0.382, 0.414]).

The regression decomposition of TPI into spatial accessibility, vulnerability, and commute burden components shows that spatial accessibility accounts for 54.2\% of TPI variance in urban village areas, compared to 61.8\% in high-end housing, suggesting that social vulnerability factors (housing type, demographic proxies) play a larger role in time poverty for lower-income populations.

%% ========== DISCUSSION ==========
\section{Discussion}
\label{sec:discussion}

\subsection{Policy Implications: Beyond Statistical Proximity}
\label{sec:policy_implications}

Our findings carry three critical implications for 15-minute city policy implementation.

\textbf{First, statistical accessibility thresholds must be supplemented with ground-truth walkability assessment.} The Accessibility Illusion Index demonstrates that network-distance-based accessibility metrics systematically overestimate experienced accessibility in informal settlement areas. A residential location may satisfy the 1,250-meter threshold to a primary school yet remain effectively inaccessible due to narrow alleyways, informal crossings, and absent pedestrian infrastructure. Policy frameworks that rely solely on statistical proximity without infrastructure quality calibration risk allocating resources to areas that already appear well-served while overlooking genuine accessibility deficits.

\textbf{Second, the 15-minute city must be disaggregated by housing type and population vulnerability.} Our finding that 31.4\% of urban village residents face time poverty, compared to 3.2\% in high-end housing, demonstrates that aggregate accessibility metrics mask profound intra-urban disparities. The 15-minute city's equity promise can only be assessed through vulnerability-weighted analyses that center the spatial experience of marginalized populations.

\textbf{Third, the urban village-peripheral accessibility gap requires targeted spatial intervention.} The LISA cluster analysis identifies the northeastern urban village axis (Baishizhou-Daxili-Taoyuan) as the primary accessibility deprivation hotspot. Our TPI results indicate that spatial interventions in this corridor---specifically pedestrian network completion, sidewalk widening, and barrier-free infrastructure retrofitting---would disproportionately benefit the most vulnerable populations.

Concrete policy recommendations include: (1) integrating street-view-based walkability audits into the 15-minute city monitoring framework; (2) prioritizing pedestrian infrastructure investment in Q4 (Accessibility Illusion) zones; (3) developing housing-type-disaggregated accessibility targets alongside aggregate district-level goals; and (4) implementing the AII as a standard metric in urban planning impact assessments for new developments.

\subsection{Methodological Contributions and Limitations}
\label{sec:method_contributions}

The integrated multi-model framework provides more robust accessibility estimates than single-model approaches: the three models show 91-93\% correlation, confirming convergent validity, while their distributional differences reveal important tail behavior that affects identification of deprivation hotspots.

The Accessibility Illusion Index represents the primary methodological innovation, operationalizing the concept of statistical accessibility-over-experienced accessibility gap into a quantifiable metric. However, several limitations warrant acknowledgment:

\begin{enumerate}
    \item \textbf{Ground-truth assessment limitations.} The LLM-vision scoring approach, while efficient and scalable, introduces evaluator subjectivity. The $r = 0.87$ inter-rater reliability, while acceptable, indicates room for improvement through structured scoring rubrics or ensemble scoring across multiple model instances. Additionally, street-view imagery captures visual features but cannot assess dynamic factors (construction activity, temporary closures, vehicle obstruction frequency) that affect real-time accessibility.
    
    \item \textbf{Housing type classification uncertainty.} The SouFun dataset (2017) provides a historical snapshot of housing prices and types; current market dynamics may have altered the price thresholds used for classification. The keyword-based urban village detection may miss some informal settlements not labeled with village-related terminology.
    
    \item \textbf{Data currency.} The POI data (2026) and housing data (2017) span different temporal snapshots, introducing potential misclassification of rapidly changing urban areas. The urban village demolition wave in Shenzhen (particularly Baishizhou, demolished 2021-2024) means some analyzed locations no longer reflect current conditions.
    
    \item \textbf{Individual-level heterogeneity.} The analysis treats housing units as the unit of analysis, aggregating residents within each housing estate. Individual-level heterogeneity (age, disability status, employment schedule, household composition) is proxied through housing-type vulnerability weights, which may not fully capture within-housing-type variation.
    
    \item \textbf{Causal inference limitations.} The correlational design cannot establish causal relationships between housing type and accessibility deprivation. The observed associations may reflect confounding by neighborhood age, land use mix, historical planning decisions, and self-selection of income-constrained households into lower-accessibility areas.
\end{enumerate}

\subsection{Integration with Existing Literature}
\label{sec:literature_integration}

Our findings extend the 2SFCA literature by demonstrating that the M2SFCA-Gaussian-Hansen framework, while robust at aggregate scales, systematically underestimates intra-urban inequality when applied to heterogeneous urban environments with informal settlements. This is consistent with \citet{dai2010spatial}'s original observation that 2SFCA variants produce biased estimates in areas with irregular service distribution, but extends this critique to the specific case of Chinese urban villages.

The time poverty findings align with \citet{pineda2020time}'s spatial time poverty framework while extending it to the Chinese megacity context. The finding that 31.4\% of urban village residents experience time poverty is comparable to \citet{li2022time}'s estimate for elderly Beijing residents (47\% time poverty for healthcare travel), suggesting that urban village populations face time poverty burdens comparable to elderly populations despite being in the primary working-age demographic.

The spatial autocorrelation results (Moran's I = 0.647) are notably higher than those reported in comparable studies \citep{vita2019spatial}, likely reflecting Shenzhen's pronounced spatial fragmentation between planned and informal urban fabrics.

%% ========== CONCLUSION ==========
\section{Conclusion}
\label{sec:conclusion}

This study developed an integrated multi-model spatial accessibility framework combining M2SFCA, Gaussian 2SFCA, and Hansen integral accessibility methods to analyze the 15-minute city performance in Shenzhen's Nanshan District. The key findings are:

\begin{enumerate}
    \item Urban village residents face a 48\% accessibility deficit (0.384 vs. 0.741) compared to high-end housing residents, despite occupying adjacent or overlapping spatial footprints.
    \item The Accessibility Illusion Index reveals that 38.6\% of urban village locations exhibit significant accessibility overestimation (AII $> 0.4$), compared to only 1.4\% in high-end housing.
    \item The Time Poverty Index shows that 31.4\% of urban village residents experience compound spatial-vulnerability time poverty, compared to 3.2\% in high-end communities.
    \item Spatial autocorrelation analysis confirms that accessibility deprivation is non-randomly clustered (Moran's I = 0.647, $p < 0.001$), with significant LISA clusters in the northeastern urban village axis.
    \item The 15-minute city's policy promise is fulfilled in only 23.4\% of the study area (Q1: True Accessibility), while 35.7\% exhibits the Accessibility Illusion (Q4) and 29.1\% faces True Deprivation (Q3).
\end{enumerate}

These findings challenge the universal applicability of the 15-minute city framework and call for a paradigm shift from statistical proximity to ground-truth pedestrian accessibility. The proposed Accessibility Illusion Index provides a replicable methodology for calibrating spatial accessibility assessments against street-level reality, applicable to any urban context where street-view imagery and network distance data are available.

Future research directions include: (1) longitudinal analysis tracking changes in AII as urban village redevelopment proceeds; (2) integration of real-time pedestrian infrastructure monitoring (IoT sensors, OpenStreetCam) for dynamic ground-truth assessment; (3) agent-based modeling to simulate the impact of targeted pedestrian infrastructure investments on AII and TPI; and (4) cross-city comparative analysis applying the AII framework to Paris, Melbourne, Portland, and other cities where 15-minute city metrics are operationalized.

\vskip 0.3in
\noindent\textbf{Acknowledgments.} This research was supported by [funding source]. We thank [acknowledgments]. Street-view imagery was processed using Anthropic Claude API. Road network data from OpenStreetMap contributors. POI data from OpenStreetMap and 大众点评.

\vskip 0.2in
\noindent\textbf{Declaration of competing interest.} The authors declare no competing interests.

%% ========== REFERENCES ==========
\bibliographystyle{apalike}
\bibliography{references}

%% ========== APPENDIX ==========
\appendix
\section{Appendix A: Model Specification Details}
\label{appendix:model_spec}

\subsection{Distance Decay Functions}

The M2SFCA employs a normalized Gaussian decay function:
\begin{equation}
f_G(d, d_0) = \frac{e^{-0.5 \cdot (d/d_0)^2} - e^{-0.5}}{1 - e^{-0.5}}
\end{equation}

The Gaussian 2SFCA uses a standard Gaussian kernel with bandwidth $\sigma = d_0 / 3$:
\begin{equation}
G(d; \sigma) = \exp\left(-\frac{d^2}{2\sigma^2}\right)
\end{equation}

\subsection{Catchment Radius Specifications}

Table \ref{tab:catchment_specs} provides the full specification of catchment radii and distance decay parameters for each service category.

\begin{threeparttable}
\caption{Catchment Radius and Distance Decay Parameters by Service Category}
\label{tab:catchment_specs}
\begin{tabular}{lccc}
\toprule
Service Category & Catchment Radius ($d_0$) & Decay Type & $\sigma$ (Gaussian model) \\
\midrule
Primary Education & \SI{1250}{\meter} & Gaussian & \SI{417}{\meter} \\
Healthcare Facilities & \SI{1250}{\meter} & Gaussian & \SI{417}{\meter} \\
Daily Commercial & \SI{1000}{\meter} & Gaussian & \SI{333}{\meter} \\
Parks/Green Space & \SI{1000}{\meter} & Gaussian & \SI{333}{\meter} \\
Public Transit & \SI{500}{\meter} & Gaussian & \SI{167}{\meter} \\
Elder Care & \SI{800}{\meter} & Gaussian & \SI{267}{\meter} \\
\bottomrule
\end{tabular}
\end{threeparttable}

\section{Appendix B: LLM-Vision Scoring Prompt Template}
\label{appendix:prompt_template}

The following structured prompt was used for LLM-vision scoring of street-view imagery:

\begin{verbatim}
You are an urban walkability expert analyzing street-view imagery 
to assess pedestrian accessibility. For each image, provide integer 
scores (0-10, half-point precision) on four dimensions:

1. Walkability Score (WS): Presence/quality of sidewalks, 
   path continuity, surface evenness, obstruction-free space.

2. Safety Index (SI): Street lighting, traffic separation, 
   crossing facilities, surveillance, personal safety.

3. Accessibility Index (AI): Wheelchair ramps, tactile paving, 
   handrails, barrier-free transitions.

4. Night Visibility Score (NVS): Estimated nighttime walkability 
   based on lighting infrastructure indicators.

Provide your response in the format:
WS: X.X [one-sentence justification]
SI: X.X [one-sentence justification]
AI: X.X [one-sentence justification]
NVS: X.X [one-sentence justification]
\end{verbatim}

\section{Appendix C: Additional Statistical Results}
\label{appendix:additional_stats}

Table \ref{tab:pairwise_tukey} presents the full pairwise Tukey HSD results for composite accessibility by housing type.

\begin{threeparttable}
\caption{Pairwise Tukey HSD Comparisons of Composite Accessibility by Housing Type}
\label{tab:pairwise_tukey}
\begin{tabular}{lccc}
\toprule
Comparison & Mean Diff. & $p$-value & Significance \\
\midrule
UV vs. AH & -0.137 & $< 0.001$ & *** \\
UV vs. CH & -0.228 & $< 0.001$ & *** \\
UV vs. HH & -0.357 & $< 0.001$ & *** \\
AH vs. CH & -0.091 & $< 0.001$ & *** \\
AH vs. HH & -0.220 & $< 0.001$ & *** \\
CH vs. HH & -0.129 & $< 0.001$ & *** \\
\bottomrule
\end{tabular}
\begin{tablenotes}
\item UV = Urban Village; AH = Affordable Housing; CH = Commodity Housing; HH = High-End Housing. *** $p < 0.001$.
\end{tablenotes}
\end{threeparttable}

\end{document}
